\documentclass[11pt]{article}
\usepackage[margin=1in]{geometry}
\usepackage{hyperref}
\usepackage{xcolor}
\usepackage{fancyvrb}
\usepackage{enumitem}
\usepackage{titlesec}
\usepackage{longtable}
\usepackage{booktabs}
\usepackage{parskip}

\definecolor{accent}{HTML}{C42130}
\definecolor{codebg}{HTML}{F5F5F7}
\definecolor{linkcol}{HTML}{1A5FB4}

\hypersetup{
  colorlinks=true,
  linkcolor=accent,
  urlcolor=linkcol,
  pdftitle={Splitting ARENA Notebooks into Self-Contained Drills},
  pdfauthor={Delta Drills}
}

\titleformat{\section}{\Large\bfseries\color{accent}}{\thesection}{0.6em}{}
\titleformat{\subsection}{\large\bfseries}{\thesubsection}{0.5em}{}

\newcommand{\colab}[1]{\href{https://colab.research.google.com/github/AkiraTheSquid/Delta-Drills/blob/main/arena-procedural-drills/#1}{Open in Colab \,$\nearrow$}}

\DefineVerbatimEnvironment{code}{Verbatim}{%
  fontsize=\small,frame=single,framerule=0pt,rulecolor=\color{codebg},%
  formatcom=\color{black},baselinestretch=0.95,xleftmargin=4pt,xrightmargin=4pt}

\setlength{\parindent}{0pt}
\setlength{\parskip}{6pt}

\title{\color{accent}\bfseries Splitting ARENA Notebooks into\\ Self-Contained Drills\\[4pt]
\large\color{black}\normalfont A worked case study: where the split was genuinely hard, and how we did it anyway}
\author{Delta Drills \quad\textbar\quad \href{https://delta-drills.vercel.app}{delta-drills.vercel.app}}
\date{June 2026}

\begin{document}
\maketitle

\section*{The claim under test}

\begin{quote}
\itshape ``While there are some existing efforts to address this issue (MiniARENA, CAMBRIA,
ARBOx3), ARENA notebooks are designed to be completed in continuous segments, making it
difficult to split and effectively run the notebooks in shorter, self-contained sessions.''
\end{quote}

The premise is correct for \emph{stock} ARENA: each notebook is a continuous narrative in
which exercise $N$ depends on definitions, imports and fixtures established in cells
$1\ldots N{-}1$. Pull any single exercise out and it will not run.

This document is the counter-example. Delta Drills decomposed ARENA Chapter~0 into
\textbf{2{,}538 standalone notebooks} --- 582 single-atom drills, 180 composite drills, and
per-question solution notebooks --- where \emph{every exercise carries its own setup} and
runs in a fresh Colab kernel with nothing executed before it. The hard part wasn't volume;
it was the topics where the continuous narrative \emph{is} the point --- the autograd engine,
the convolution-as-strided-view, the DCGAN decoder. Those are the examples below.

\section{How a split exercise is made self-contained}

The mechanism is the same in every drill, and it is what defeats ``continuous segments'':
\textbf{each notebook re-declares the minimal scaffolding it needs, inline, inside its own
test harness.} It never imports the rest of the chapter.

A drill notebook is a fixed 6-part skeleton:

\begin{enumerate}[leftmargin=1.4em]
  \item \textbf{Setup cell} --- imports + \verb|t.manual_seed(0); np.random.seed(0)|. Deterministic, self-contained.
  \item \textbf{Auth cell} --- \verb|DD_TOKEN|, \verb|DD_ATOM_IDS|, \verb|DD_SUBTOPICS|, backend URL.
  \item \textbf{``How these atoms compose''} --- the prose that replaces the missing notebook context.
  \item \textbf{Exercise stub} --- a single function with \verb|raise NotImplementedError|.
  \item \textbf{Inline test harness} --- re-declares the local \verb|Recipe| / \verb|MiniTensor| /
        \verb|back_funcs| etc.\ that the real ARENA notebook would have built over many cells, then asserts correctness.
  \item \textbf{Completion beacon} --- POSTs to \verb|/api/practice/arena-rating|, updating BKT mastery for every atom the drill exercises.
\end{enumerate}

The crucial line is in the grader (mirrored inside every notebook's test): seed \emph{then}
run per-exercise setup, so the learner's function executes against a known fixed state with
no dependence on prior cells.

\begin{code}
np.random.seed(0)                       # known state
exec(case["setup_code"], globals())     # ONLY this exercise's fixtures
actual = eval(case["call"], globals())  # run solve() in isolation
\end{code}

\section{Examples where the split was genuinely hard --- and we did it}

The autograd chapter (ARENA \texttt{[0.4] Backprop}) is the hardest case: it builds a working
mini-autograd engine where \emph{everything depends on everything} --- a \texttt{Tensor}
wrapper, a \texttt{Recipe} of how each value was computed, a topological sort, and a reverse
loop that dispatches per-argument backward functions. You cannot test the reverse loop without
the graph; you cannot build the graph without the wrapper. Stock ARENA does this across one
long notebook. We split it into atoms that each \emph{rebuild just enough of the engine to
stand alone.}

\subsection{The full reverse pass, in one isolated cell}
\textbf{cx25 --- seed reverse pass with \texttt{ones\_like}; pop-and-dispatch loop.} \;\colab{composites/part4/025-end-grad-seeded-backprop-loop.ipynb}

This is the heart of the whole chapter. The notebook re-declares a \verb|Recipe| dataclass, a
\verb|MiniTensor| class, and two \verb|multiply_back| functions \emph{inside the test}, so the
reverse-mode loop can be exercised with zero upstream cells:

\begin{code}
def cx25_backprop(end_node, end_grad, sorted_graph, back_funcs):
    if end_grad is None:
        seed = t.ones_like(end_node.array)          # d(end.sum())/d(end)
    else:
        assert end_grad.array.shape == end_node.array.shape
        seed = end_grad.array
    grads = {id(end_node): seed}
    for node in sorted_graph:                        # reverse-pass driver
        grad_out = grads.pop(id(node), None)
        if grad_out is None: continue
        if node.recipe is None:                      # leaf -> .grad
            node.grad = grad_out if node.grad is None else node.grad + grad_out
            continue
        for argnum, parent in node.recipe.parents.items():
            back = back_funcs[(node.recipe.func, argnum)]
            gp = back(grad_out, node.array, *node.recipe.args, **node.recipe.kwargs)
            grads[id(parent)] = grads.get(id(parent), 0) + gp
\end{code}

Why it is hard to split, and why this split is good: the seed step and the loop step are
\emph{inseparable} in ARENA --- the loop can't start without a seed in \texttt{grads}, and the
seed has no purpose without the loop. Rather than fake that by stubbing the engine, the drill
keeps both atoms in one function and supplies a tiny but \emph{real} graph in the test. The
learner writes the exact code they'd write in ARENA.

\subsection{Topological sort with cycle detection}
\textbf{cx19 --- three-colour DFS: deps-first order \emph{and} \texttt{ValueError} on a cycle.} \;\colab{composites/part4/019-dfs-toposort-with-temp-set-cycle-detection.ipynb}

A single DFS walk carries two responsibilities --- two atoms --- that ARENA never separates:
the \texttt{perm}/\texttt{temp} colour sets (cycle detection) and the post-order append
(deps-first ordering). The test feeds a linear chain, a diamond DAG (node must appear once),
and a self-loop (must raise), so passing proves both atoms compose correctly.

\begin{code}
def cx19_topo_sort_with_cycle_check(root, get_children):
    result, perm, temp = [], set(), set()
    def visit(node):
        nid = id(node)
        if nid in perm: return                     # finished subtree, fine
        if nid in temp: raise ValueError(f'cycle at {node!r}')  # back-edge
        temp.add(nid)
        for child in get_children(node): visit(child)
        temp.remove(nid); perm.add(nid)
        result.append(node)                         # deps-first append
    visit(root); return result
\end{code}

\subsection{Wiring a brand-new op into the engine}
\textbf{cx1 --- register a forward op so it builds a \texttt{Recipe} and back-propagates.} \;\colab{composites/part4/001-wire-new-op-into-autograd.ipynb}

\textbf{cx30 --- \texttt{unbroadcast} inverts NumPy broadcasting in the backward pass.} \;\colab{composites/part4/030-unbroadcast-inverts-broadcast.ipynb}

These two are the ``glue'' atoms that, in stock ARENA, are impossible to test until the entire
engine exists. Splitting them out required re-creating a minimal differentiable-op registry in
the test cell --- which is exactly the elegant part: the drill proves you understand the
contract (\texttt{forward} builds a recipe; \texttt{back} replays it; gradients un-broadcast
back to parameter shape) without the 400-line notebook around it.

\section{Elegant splits: hard core problem, clean two-line drill}

These are cases where the underlying ARENA exercise is conceptually heavy, but the atomised
drill is short and sharp --- the split \emph{distils} the difficulty instead of diluting it.

\subsection{Conv2d forward as an \texttt{as\_strided} view + einsum}
\textbf{cx2 --- the whole convolution in two lines.} \;\colab{composites/part2/002-as-strided-then-einsum-conv-forward.ipynb}

ARENA's convolution exercise is one of the chapter's most feared. The atomised version reduces
it to its two load-bearing ideas --- a strided patch view and a single einsum contraction:

\begin{code}
def cx2_conv2d_via_einsum(x, w):
    B, C_in, H, W = x.shape
    C_out, _, KH, KW = w.shape
    H_out, W_out = H - KH + 1, W - KW + 1
    sB, sC, sH, sW = x.stride()
    patches = x.as_strided(size=(B, C_in, H_out, W_out, KH, KW),
                           stride=(sB, sC, sH, sW, sH, sW))   # spatial stride twice
    return einops.einsum(patches, w, 'b c h w kh kw, o c kh kw -> b o h w')
\end{code}

The elegance is in the stride tuple: the spatial strides appear \emph{twice} because the same
2-D step walks both the patch origin and the inside of each patch. The drill isolates exactly
that insight. (Companion drill cx1 \,\colab{composites/part2/001-as-strided-view-sized-by-conv-output-shape.ipynb}
covers deriving the output shape that sizes the view.)

\subsection{Batched ray/triangle solve}
\textbf{cx2 --- stack columns, then one batched \texttt{linalg.solve}.} \;\colab{composites/part1/002-stack-then-batched-solve.ipynb}

\begin{code}
def cx2_solve_from_columns(cols, b):
    A = t.stack(cols, dim=-1)     # NEW trailing axis -> (K, n, n)
    return t.linalg.solve(A, b)   # K LU factorisations fused into one BLAS call
\end{code}

Two atoms, two lines. The drill hinges on the \texttt{stack} vs \texttt{cat} choice ---
\texttt{cat(dim=-1)} would silently give \texttt{(K, n*n)} and either crash or, worse, return
garbage. (Companion cx4 \,\colab{composites/part1/004-patch-singular-then-batched-solve.ipynb}
adds the singular-slice patch that real ray tracing needs.)

\subsection{The DCGAN decoder, as one \texttt{nn.Module}}
\textbf{cx4 --- latent $z \to$ Linear projection $\to$ reshape $\to$ ConvT/BN/ReLU upsample.} \;\colab{composites/part5/004-latent-projection-then-convt-block.ipynb}

\begin{code}
class Decoder(nn.Module):
    def __init__(self, latent_dim, base_channels, base_hw):
        super().__init__()
        self.base_c, self.base_hw = base_channels, base_hw
        self.proj = nn.Linear(latent_dim, base_channels * base_hw * base_hw)  # bare, no act
        self.block = nn.Sequential(
            nn.ConvTranspose2d(base_channels, base_channels // 2, 4, 2, 1, bias=False),
            nn.BatchNorm2d(base_channels // 2),
            nn.ReLU(inplace=True))
\end{code}

A generative-model component that ordinarily only makes sense inside a full training notebook
becomes a self-contained module-construction drill. (Companion cx1
\,\colab{composites/part5/001-dcgan-g-block-as-nn-module-subclass.ipynb} drills the
ConvT/BN/activation block on its own.)

\subsection{Optimizer internals}
\textbf{cx2 --- momentum SGD \texttt{\_\_init\_\_}: materialise params \emph{and} per-param velocity buffer.} \;\colab{composites/part3/002-sgd-init-with-velocity-buffers.ipynb}

\textbf{cx4 --- update the velocity buffer in place via \texttt{copy\_}.} \;\colab{composites/part3/004-velocity-buffer-update-via-copy.ipynb}

The optimizer loop is another ARENA exercise that only ``works'' once the whole training script
exists. These drills isolate the two subtle bits --- buffer allocation at init, and in-place
update --- each runnable against a couple of throwaway tensors.

\section{Full example index}

Every link opens a live Colab notebook from the \texttt{AkiraTheSquid/Delta-Drills} repo
(\texttt{main} branch). Signed-in students get their own fork via the account GitHub-username
field, so Colab ``Save a copy in GitHub'' persists their work.

\renewcommand{\arraystretch}{1.3}
\begin{longtable}{@{}p{2.1cm}p{8.1cm}p{2.4cm}@{}}
\toprule
\textbf{Domain} & \textbf{Drill} & \textbf{Link} \\
\midrule
\endhead
Autograd & cx25 seed + pop-and-dispatch reverse pass & \colab{composites/part4/025-end-grad-seeded-backprop-loop.ipynb} \\
Autograd & cx19 DFS toposort + cycle detection & \colab{composites/part4/019-dfs-toposort-with-temp-set-cycle-detection.ipynb} \\
Autograd & cx1 wire a new op into autograd & \colab{composites/part4/001-wire-new-op-into-autograd.ipynb} \\
Autograd & cx30 unbroadcast inverts broadcast & \colab{composites/part4/030-unbroadcast-inverts-broadcast.ipynb} \\
CNN & cx2 Conv2d via as\_strided + einsum & \colab{composites/part2/002-as-strided-then-einsum-conv-forward.ipynb} \\
CNN & cx1 as\_strided view sized by conv output shape & \colab{composites/part2/001-as-strided-view-sized-by-conv-output-shape.ipynb} \\
Ray tracing & cx2 stack columns then batched solve & \colab{composites/part1/002-stack-then-batched-solve.ipynb} \\
Ray tracing & cx4 patch singular slices then solve & \colab{composites/part1/004-patch-singular-then-batched-solve.ipynb} \\
Optimizers & cx2 momentum SGD init + velocity buffers & \colab{composites/part3/002-sgd-init-with-velocity-buffers.ipynb} \\
Optimizers & cx4 velocity buffer update via copy\_ & \colab{composites/part3/004-velocity-buffer-update-via-copy.ipynb} \\
VAE/GAN & cx4 latent projection then ConvT block & \colab{composites/part5/004-latent-projection-then-convt-block.ipynb} \\
VAE/GAN & cx1 DCGAN G block as nn.Module subclass & \colab{composites/part5/001-dcgan-g-block-as-nn-module-subclass.ipynb} \\
Einops & cx2 outer-product two 1-D grids via repeat & \colab{composites/part0/002-outer-product-grids-via-repeat.ipynb} \\
\bottomrule
\end{longtable}

The full library is browsable in-app at \href{https://delta-drills.vercel.app}{delta-drills.vercel.app}
(Practice tab) and in the repo under \texttt{arena-procedural-drills/}.

\section{Method: how the splits were extracted}

The split was a pipeline, not a manual port. Each stage \emph{recycled the artifact built for
testing the previous stage} --- which is what made it cheap to scale to thousands of notebooks.

\subsection*{1. Atomise the curriculum}
ARENA's exercises were decomposed into a concept graph of \textbf{207 drillable atoms} (e.g.\
\texttt{as-strided-windowing}, \texttt{dfs-three-set-toposort}, \texttt{end-grad-default-ones-like}),
with prerequisite and encompassing edges. Each atom is a single testable skill. Composite
drills then exercise 2--3 atoms \emph{together}, because some atoms (the seed and the reverse
loop) are inseparable in real code.

\subsection*{2. Author standalone scaffolding per atom}
For each atom a notebook was generated carrying its own imports, deterministic seed, a function
stub, and --- critically --- an \textbf{inline test harness that rebuilds the minimal engine}
(\texttt{Recipe}, \texttt{MiniTensor}, \texttt{back\_funcs}, helper graph nodes) the exercise
needs. This is the step that breaks ``continuous segments'': the context that ARENA accumulates
across cells is re-created locally, in miniature, in one cell.

\subsection*{3. Recycle the grader as the validator}
The same harness that scores a student's submission at runtime
(\texttt{code\_runner.run\_function\_tests}: seed $\to$ run \texttt{setup\_code} $\to$ call
\texttt{solve()}) was reused at \emph{build time} to validate every authored solution. A drill
ships only if the reference solution passes its own test under the live grader. The test
fixture, the validation oracle, and the production grader are one object --- not three.

This recycling surfaced a real, load-bearing quirk: the grader seeds the RNG \emph{then} runs
\texttt{setup\_code} (which consumes the stream), so a correct \texttt{solve()} must
\emph{reseed internally} before drawing. Because the validator \emph{is} the grader, every
authored solution was forced to obey the same contract students face --- no drift between
``passes in authoring'' and ``passes in the app''.

\subsection*{4. Build solution + hint siblings from the same source}
For each problem notebook a \texttt{.solution.ipynb} sibling was generated by lifting the code
out of the \verb|<details><summary>Show solution</summary>| block into the
\texttt{NotImplementedError} stub (761/762 drills), then re-running the grader to confirm the
solution notebook itself executes clean. Per-drill nudge hints were authored and validated the
same way. The ``Show answer'' / ``Show hint'' buttons in the app point at these siblings,
routed to the student's Colab fork.

\subsection*{5. Wire completion back to mastery}
Each notebook ends with a beacon POSTing to the backend, which updates Bayesian Knowledge
Tracing mastery for every atom the drill exercises and unlocks downstream atoms through the
prerequisite graph. The split isn't just \emph{runnable} in a short session --- it
\emph{closes the loop} on the adaptive engine.

\section*{Verdict}

The claim's premise holds for stock ARENA; its implied limitation does not hold against this
system. Delta Drills shipped the self-contained-session split the claim treats as the open
problem --- including on the chapters (autograd, convolution, DCGAN) where the continuous
narrative is hardest to break --- with reproducible per-exercise setup, isolated grading reused
as the build-time validator, and mastery wired back through the concept graph. The honest
framing: ARENA is hard to split \emph{by default}; we solved it with atomisation plus
self-contained scaffolding, not by wishing the dependency away.

\end{document}
